\documentclass[11pt,a4paper]{article}

% ============================================================
%  Buku Panduan SeBatt (Secure Battery Identity Chip)
%  Gaya mengikuti buku "Implementasi CORDIC dan Transformasi Clarke-Park"
%  (POWEL/Project/fpga-cordic). Kompilasi dari folder ini:
%    latexmk -pdf buku-panduan-sebatt.tex
%  Status: kerangka. Setiap \isi{...} adalah catatan isi yang masih harus
%  ditulis, dan \isian{...} adalah angka yang menunggu hasil uji.
% ============================================================

% ============================================================
%  PAKET
% ============================================================
\usepackage[top=3cm, bottom=3cm, left=3cm, right=2.5cm]{geometry}
\usepackage[utf8]{inputenc}
\usepackage[T1]{fontenc}
\usepackage[indonesian]{babel}
\usepackage{amsmath, amssymb}
\usepackage{graphicx}
\usepackage{fancyhdr}
\usepackage{listings}
\usepackage{xcolor}
\usepackage{enumitem}
\usepackage{booktabs}
\usepackage{tabularx}
\usepackage{hyperref}
\usepackage[nobottomtitles*]{titlesec}
\usepackage[default]{lato}
\usepackage{indentfirst}
\usepackage{float}
\usepackage{caption}
\usepackage{tikz}
\usetikzlibrary{positioning, calc, fit, backgrounds}

\captionsetup{font=small, labelfont=bf, skip=5pt}
\setlength{\intextsep}{8pt plus 2pt minus 2pt}
\setlength{\textfloatsep}{8pt plus 2pt minus 2pt}

% ============================================================
%  WARNA
% ============================================================
\definecolor{headinggray}{HTML}{4A5560}
\definecolor{accentred}{HTML}{D32F2F}
\definecolor{rule}{HTML}{C7CDD2}
\definecolor{covernavy}{HTML}{1F2F4A}
\definecolor{covergold}{HTML}{B8923A}
\definecolor{covernavylight}{HTML}{3D5A80}
\definecolor{covergray}{HTML}{6B7280}

\setlist{itemsep=3pt, topsep=4pt}
% Ruang tambahan agar baris dengan nama modul \texttt panjang tidak menonjol keluar margin
\setlength{\emergencystretch}{2em}
% Gambar isi di images/, gambar teknis bersama dari docs/gambar/
\graphicspath{{images/}{../gambar/}}
\renewcommand{\figurename}{Gambar}
\renewcommand{\tablename}{Tabel}
\renewcommand{\contentsname}{Daftar Isi}
\renewcommand{\refname}{Daftar Pustaka}
\renewcommand{\lstlistingname}{Kode}

% Isian yang menunggu hasil uji: \isian{keterangan}
\newcommand{\isian}[1]{\textcolor{accentred}{[#1]}}

% Catatan kerangka (isi yang belum ditulis): \isi{keterangan}
\newcommand{\isi}[1]{%
  \par\noindent\fcolorbox{rule}{gray!6}{\parbox{\dimexpr\linewidth-2\fboxsep-2\fboxrule}{%
    \small\itshape\color{gray!70!black}[Isi] #1}}\par\medskip}

% Kotak pengganti gambar yang belum ada: \gambarplaceholder{keterangan}
\newcommand{\gambarplaceholder}[1]{%
  \begin{center}
    \fbox{\parbox[c][2.6cm][c]{0.8\textwidth}{\centering\small\itshape\color{gray}
      [GAMBAR BELUM ADA]\\[4pt] #1}}
  \end{center}
}

% Gambar dari berkas: \gambarberkas{nama.png}{lebar}{keterangan}{label}
% Selama berkasnya belum ada, yang tampil kotak berisi nama berkas yang harus disiapkan.
\newcommand{\gambarberkas}[4]{%
  \begin{figure}[H]
    \centering
    \IfFileExists{images/#1}{\includegraphics[width=#2\textwidth]{#1}}{%
    \IfFileExists{../gambar/#1}{\includegraphics[width=#2\textwidth]{#1}}{%
      \fbox{\parbox[c][2.6cm][c]{0.8\textwidth}{\centering\small\color{gray}
        [GAMBAR BELUM ADA]\\[4pt] \texttt{images/#1}}}}}
    \caption{#3}
    \label{#4}
  \end{figure}
}

% ============================================================
%  HEADER DAN FOOTER
% ============================================================
\setlength{\headheight}{14pt}
% \rightmark memberi bab pertama pada halaman, sehingga header cocok dengan judul di atasnya
\renewcommand{\sectionmark}[1]{\markboth{\thesection\quad #1}{\thesection\quad #1}}
\renewcommand{\subsectionmark}[1]{}
\renewcommand{\headrule}{{\color{rule}\hrule width\headwidth height\headrulewidth}}
\newcommand{\nomorhalaman}{%
  \tikz[baseline=(n.base)]\node[fill=covernavy, text=white, rounded corners=2pt, inner xsep=5pt,
    inner ysep=3pt, font=\small\bfseries] (n) {\thepage};}
\fancypagestyle{isi}{%
  \fancyhf{}
  \fancyhead[L]{%
    \begin{tikzpicture}[remember picture, overlay]
      \fill[covernavy] (current page.north west) rectangle
        ([xshift=0.9cm, yshift=-2.6cm]current page.north west);
    \end{tikzpicture}%
    \small\color{headinggray!70}SeBatt: Secure Battery Identity Chip pada FPGA}
  \fancyhead[R]{\small\bfseries\color{covernavy}%
    \tikz[baseline=-0.6ex]\fill[covergold] (0,0) rectangle (0.16,0.16);\hspace{0.5em}\rightmark}
  \fancyfoot[C]{\nomorhalaman}
  \renewcommand{\headrulewidth}{0.4pt}
  \renewcommand{\footrulewidth}{0pt}}
\fancypagestyle{plain}{%
  \fancyhf{}
  \fancyfoot[C]{\small\color{headinggray}\thepage}
  \renewcommand{\headrulewidth}{0pt}
  \renewcommand{\footrulewidth}{0pt}}
\pagestyle{isi}

% ============================================================
%  GAYA LISTING KODE (Verilog dan terminal)
% ============================================================
\lstset{
  language=Verilog,
  basicstyle=\ttfamily\footnotesize,
  keywordstyle=\bfseries\color{headinggray},
  commentstyle=\itshape\color{gray!70!black},
  stringstyle=\color{accentred!80!black},
  breaklines=true,
  breakindent=1.5em,
  frame=single,
  rulecolor=\color{rule},
  columns=fixed,
  basewidth=0.52em,
  keepspaces=true,
  backgroundcolor=\color{gray!5},
  showstringspaces=false,
  tabsize=2,
}
\lstdefinestyle{terminal}{language={}, basicstyle=\ttfamily\small}

% ============================================================
%  HYPERLINK DAN METADATA PDF
% ============================================================
\hypersetup{
  colorlinks=true,
  linkcolor=black,
  citecolor=covernavylight,
  urlcolor=blue,
  pdftitle={SeBatt: Secure Battery Identity Chip pada FPGA},
  pdfauthor={Guntur Sulistyo, Muhammad Budi Setiawan, Muhammad Muqtada Alhaddad},
}

\newcommand{\tanggalindo}{%
  \number\day{} %
  \ifcase\month\or Januari\or Februari\or Maret\or April\or Mei\or Juni\or
    Juli\or Agustus\or September\or Oktober\or November\or Desember\fi
  \space \number\year}

% ============================================================
%  FORMAT JUDUL SECTION
% ============================================================
\titleformat{\section}
  {\Large\bfseries\color{covernavy}}
  {\thesection}{0.6em}{}
  [{\color{covergold}\titlerule[1.2pt]}]
\titleformat{\subsection}{\large\bfseries\color{covernavy}}{\thesubsection}{0.6em}{}
\titleformat{\subsubsection}{\normalsize\bfseries\color{covernavylight}}{\thesubsubsection}{0.6em}{}

\titlespacing*{\section}{0pt}{1.6em}{0.8em}
\titlespacing*{\subsection}{0pt}{1.2em}{0.5em}
\titlespacing*{\subsubsection}{0pt}{1em}{0.4em}

\begin{document}

% ============================================================
%  SAMPUL
% ============================================================
\begin{titlepage}
  \begin{tikzpicture}[remember picture, overlay, >=stealth]
    % batang tepi kiri
    \fill[covernavy] (current page.north west) rectangle
      ([xshift=0.9cm]current page.south west);
    % judul
    \node[anchor=north west, text width=14.5cm, align=left, text=covernavy]
      at ([xshift=3cm, yshift=-5cm]current page.north west) {%
      {\fontsize{30}{38}\selectfont\bfseries\spaceskip=0.32em\xspaceskip=0.32em
        SeBatt: Secure\\Battery Identity Chip\\pada FPGA\par}};
    \fill[covergold] ([xshift=3cm, yshift=-11.4cm]current page.north west) rectangle
      ([xshift=5.5cm, yshift=-11.55cm]current page.north west);
    \node[anchor=north west, text width=14.5cm, align=left, text=covergray]
      at ([xshift=3cm, yshift=-12cm]current page.north west) {%
      {\large Riwayat Tegangan Baterai Bertanda Tangan HMAC-SHA-256,\\
        Perakitan, dan Penggunaan\par}};
    % ilustrasi: rantai catatan bertanda tangan, satu catatan disorot
    \begin{scope}[shift={([xshift=-9.6cm, yshift=4.2cm]current page.south east)}]
      \foreach \i/\s in {0/41, 1/42, 2/43, 3/44} {
        \draw[covernavy!35, line width=0.6pt, rounded corners=2pt]
          (\i*2.2, 0) rectangle ++(1.7, 1.15);
        \node[font=\scriptsize\ttfamily, text=covernavy!60, anchor=north west]
          at (\i*2.2+0.12, 1.05) {seq \s};
        \draw[covernavy!25, line width=0.5pt] (\i*2.2+0.15, 0.45) -- ++(1.4, 0);
        \draw[covernavy!25, line width=0.5pt] (\i*2.2+0.15, 0.25) -- ++(1.0, 0);
      }
      \foreach \i in {0, 1, 2}
        \draw[->, covernavy!35, line width=0.6pt] (\i*2.2+1.7, 0.575) -- ++(0.5, 0);
      \draw[covergold, line width=1.4pt, rounded corners=2pt] (3*2.2, 0) rectangle ++(1.7, 1.15);
      \fill[covergold] (3*2.2+1.45, 0.25) circle (0.07);
    \end{scope}
    % penulis
    \node[anchor=south west, align=left, text width=14cm, text=covernavy]
      at ([xshift=3cm, yshift=3cm]current page.south west) {%
      {\large\bfseries Guntur Sulistyo\par}
      \vspace{0.15cm}
      {\large\bfseries Muhammad Budi Setiawan\par}
      \vspace{0.15cm}
      {\large\bfseries Muhammad Muqtada Alhaddad\par}
      \vspace{0.3cm}
      {\normalsize\color{covergray} Tim NaS, Universitas Gadjah Mada\par}};
  \end{tikzpicture}
\end{titlepage}

\pagenumbering{roman}
\pagestyle{plain}

% ============================================================
%  PENGANTAR
% ============================================================
\section*{Pengantar}
\addcontentsline{toc}{section}{Pengantar}

Puji syukur penulis panjatkan ke hadirat Tuhan Yang Maha Esa, karena atas rahmat-Nya buku
\textit{SeBatt: Secure Battery Identity Chip pada FPGA} ini dapat diselesaikan.

Baterai motor listrik pada jaringan stasiun penukaran berpindah tangan setiap kali ditukar.
Riwayat pemakaiannya, misalnya pernah terkuras dalam, menentukan nilai sisa dan klaim garansi,
tetapi riwayat tersebut dapat diubah atau di-reset tanpa diketahui stasiun. SeBatt menandatangani
setiap sampel tegangan dengan HMAC-SHA-256 beserta nomor urutnya, sehingga stasiun dapat
memeriksa apakah riwayat sebuah baterai masih utuh.

Buku ini membahas SeBatt mulai dari konsep dasar dan protokolnya, rancangan modul Verilog, dan
perakitan rangkaian pembacaan tegangan, sampai penggunaan stasiun dan hasil pengujian di
perangkat keras. Purwarupa yang dibahas berjalan sebagai RTL pada FPGA AMD Kria KV260, membaca
baterai LiPo melalui ADC ADS1115, dan diperiksa oleh program stasiun berbasis Python. Kode
lengkap semua modul dimuat pada lampiran, sehingga pembaca dapat mengikuti dan mengulangi setiap
langkah.

Buku ini ditujukan bagi mahasiswa, peneliti, dan pengembang di bidang sistem digital dan sistem
baterai yang ingin menjalankan ulang atau melanjutkan purwarupa SeBatt. Pembaca diharapkan sudah
mengenal dasar Verilog dan Python, tetapi tidak harus sudah memahami HMAC maupun protokol
stasiun penukaran.

Purwarupa SeBatt dikembangkan dalam Hackathon Deeptech 2026: \textit{Designing Chips for
Real-World Impact} yang diselenggarakan oleh ICDeC ITB. Penulis mengucapkan terima kasih kepada
penyelenggara, para mentor, serta semua pihak yang telah membantu.

Penulis berharap buku ini bermanfaat bagi pembaca. Kritik dan saran untuk perbaikan buku ini
sangat penulis harapkan.

\vspace{1cm}
\begin{flushright}
  Bandung, \tanggalindo\\[1.5cm]
  Tim NaS
\end{flushright}

\clearpage

% ============================================================
%  DAFTAR ISI
% ============================================================
\tableofcontents
\clearpage
\pagenumbering{arabic}
\pagestyle{isi}

% ============================================================
%  1. PENDAHULUAN
% ============================================================
\section{Pendahuluan}
\label{sec:pendahuluan}

Jumlah kendaraan listrik di Indonesia mencapai 541.678 unit pada Agustus 2026
\cite{okezone}. Sebagian sepeda motor listrik memakai baterai yang ditukar di stasiun penukaran.
SWAP Energi mengoperasikan lebih dari 1.500 stasiun \cite{swapid}, Electrum lebih dari 400
stasiun \cite{electrum}, dan V-Green lebih dari 2.000 titik penukaran \cite{detikvgreen}.
Pertumbuhan ini juga diperhitungkan dalam perencanaan ketenagalistrikan nasional. Rencana Umum
Ketenagalistrikan Nasional (RUKN) memproyeksikan kebutuhan listrik naik dari sekitar 539~TWh
pada 2025 menjadi sekitar 1.813~TWh pada 2060, dengan kendaraan bermotor listrik menyumbang
sekitar 11\% dari kebutuhan tahun 2060 \cite{rukn}. Rencana Usaha Penyediaan Tenaga Listrik
(RUPTL) PT PLN (Persero) 2025--2034 memasukkan kendaraan listrik sebagai salah satu pendorong
pertumbuhan kebutuhan listrik \cite{paparanruptl} dan mencakup penyediaan stasiun penukaran
baterai kendaraan listrik umum (SPBKLU) untuk kendaraan roda dua \cite{databoksruptl}.

Pada model penukaran, baterai dimiliki operator, sehingga satu baterai berpindah dari satu pengendara ke
pengendara lain setiap kali ditukar. Harga satu baterai mencapai Rp4,55 juta \cite{detikvgreen}.
Nilai sisa dan klaim garansinya bergantung pada riwayat pemakaian, misalnya berapa kali baterai
pernah terkuras di bawah batas aman.

\textit{Gauge} baterai yang beredar sudah merekam riwayat pemakaian. Sebagai contoh, BQ40Z50 dari
Texas Instruments menyimpan \textit{Lifetime Data} dan \textit{Black Box Recorder}, tetapi riwayat
tersebut hanya dilindungi mode akses \texttt{SEALED} dan \texttt{UNSEALED} serta dapat dihapus
melalui perintah \textit{Lifetime Data Reset} \cite{bq40z50}. Autentikasi pada \textit{gauge}
tersebut hanya membuktikan identitas baterai, sedangkan data ukurnya tidak ikut ditandatangani.
Akibatnya, stasiun yang membaca riwayat tidak dapat membedakan riwayat asli dari riwayat yang
telah diubah atau di-\textit{reset}.

SeBatt menandatangani setiap sampel tegangan di dalam rangkaian yang sama dengan pembacanya.
Setiap sampel diberi nomor urut yang hanya bertambah, lalu ditandatangani dengan HMAC-SHA-256
bersama identitas baterai. Saat baterai ditukar, stasiun meminta seluruh riwayat dengan nonce
baru dan menghitung ulang setiap tanda tangan, sehingga catatan yang diubah, dihapus, atau
diputar balik akan ditolak. Buku ini memandu pembuatan dan penggunaan SeBatt sampai terbukti
bekerja di FPGA. Setiap bagian dikerjakan dengan urutan yang sama, yaitu konsep, protokol, modul
Verilog, verifikasi dengan model Python dan \textit{testbench} cocotb, lalu uji di board.

{\sloppy
Modul yang dibahas adalah \texttt{auth\_ctrl}, \texttt{hmac\_sha256}, \texttt{sha256\_core},
\texttt{uart\_if}, \texttt{ads1115\_ctrl}, dan \texttt{i2c\_master}, yang dirangkai oleh
\texttt{sebatt\_top} dan \texttt{kv260\_top}. Semua modul ditulis dalam Verilog standar tanpa IP
vendor dan tanpa prosesor lunak. Rancangan dijalankan pada board AMD Kria KV260 dengan
\textit{clock} 100~MHz dari PS. Tegangan baterai LiPo 3S dibaca melalui modul ADC ADS1115,
sedangkan stasiun berjalan sebagai program Python di laptop yang terhubung melalui adaptor UART
CP2102. Purwarupa ini masih memakai kunci tetap, penghitung nomor urut yang kembali ke nol saat
daya dicabut, dan ADC berupa modul terpisah. Batasan tersebut dibahas pada Bab~\ref{sec:masalah}.
\par}

Pembaca diasumsikan sudah mengenal dasar Verilog, seperti \texttt{module}, \texttt{reg}, dan
blok \texttt{always}, pernah memakai Vivado, serta dapat menjalankan program Python dari
terminal. Pembaca tidak harus sudah memahami HMAC, nonce, atau protokol stasiun penukaran, karena
konsep tersebut dijelaskan pada Bab~\ref{sec:konsep}.

Gambar~\ref{fig:petabuku} memperlihatkan hubungan antarbab. Bab~\ref{sec:konsep} menjelaskan
tanda tangan HMAC, nonce, nomor urut, dan pembacaan tegangan. Konsep tersebut dipakai
Bab~\ref{sec:arsitektur} untuk menyusun format frame, alur penukaran, dan kode status.
Bab~\ref{sec:rtl} mewujudkan protokol tersebut menjadi modul Verilog beserta verifikasinya.
Bab~\ref{sec:rakit} merakit rangkaian pembacaan tegangan dan sambungannya ke board.
Bab~\ref{sec:panduan} memandu penggunaan stasiun dalam skenario normal maupun serangan. Hasil
pengujiannya di board dimuat pada Bab~\ref{sec:hasil}. Pemecahan masalah, peringatan, dan batasan
purwarupa dimuat pada Bab~\ref{sec:masalah}.

\begin{figure}[H]
  \centering
  \resizebox{\textwidth}{!}{%
  \begin{tikzpicture}[>=stealth, font=\footnotesize,
      blok/.style={draw, rounded corners, fill=white, minimum width=2.4cm,
        minimum height=0.7cm, align=center}]
    \node[blok, minimum height=2.3cm] (ks) at (0,0) {HMAC, nonce,\\nomor urut,\\pembagi\\tegangan};
    \node[blok, minimum height=2.3cm] (pr) at (3.2,0) {Frame dan\\perintah,\\kode status};
    \node[blok] (au) at (6.6,0.85) {\texttt{auth\_ctrl}};
    \node[blok] (hm) at (6.6,0) {\texttt{hmac\_sha256}};
    \node[blok] (ad) at (6.6,-0.85) {\texttt{ads1115\_ctrl}};
    \node[blok] (rk) at (10.0,0.45) {Rangkaian\\LiPo + ADC};
    \node[blok] (pm) at (10.0,-0.6) {Pmod J2};
    \node[blok] (st) at (13.4,0.45) {Stasiun dan\\\textit{dashboard}};
    \node[blok] (sk) at (13.4,-0.6) {Skenario uji};
    \begin{scope}[on background layer]
      \node[draw=gray!60, fill=gray!6, rounded corners, inner sep=7pt, fit=(ks),
        label={[font=\small]above:Bab~\ref{sec:konsep}}] (b2) {};
      \node[draw=gray!60, fill=gray!6, rounded corners, inner sep=7pt, fit=(pr),
        label={[font=\small]above:Bab~\ref{sec:arsitektur}}] (b3) {};
      \node[draw=gray!60, fill=gray!6, rounded corners, inner sep=7pt, fit=(au)(hm)(ad),
        label={[font=\small]above:Bab~\ref{sec:rtl}}] (b4) {};
      \node[draw=gray!60, fill=gray!6, rounded corners, inner sep=7pt, fit=(rk)(pm),
        label={[font=\small]above:Bab~\ref{sec:rakit}}] (b5) {};
      \node[draw=gray!60, fill=gray!6, rounded corners, inner sep=7pt, fit=(st)(sk),
        label={[font=\small]above:Bab~\ref{sec:panduan}--\ref{sec:hasil}}] (b6) {};
    \end{scope}
    \draw[->] (b2.east) -- (b3.west);
    \draw[->] (b3.east) -- (b4.west);
    \draw[->] (b4.east) -- (b5.west);
    \draw[->] (b5.east) -- (b6.west);
  \end{tikzpicture}}
  \caption{Peta buku dan bagian yang dibangun pada setiap bab.}
  \label{fig:petabuku}
\end{figure}

% ============================================================
%  2. KONSEP DASAR
% ============================================================
\section{Konsep Dasar}
\label{sec:konsep}

Bab ini menjelaskan empat konsep yang dipakai di seluruh buku, yaitu tanda tangan HMAC-SHA-256,
nonce, nomor urut monoton, dan pembacaan tegangan baterai melalui pembagi tegangan. Setiap
konsep langsung dikaitkan dengan perannya di SeBatt dan dengan kode status yang dilaporkan
stasiun. Istilah lain yang muncul di buku ini dirangkum pada Lampiran~\ref{app:istilah}.

\subsection{Tanda Tangan HMAC-SHA-256}
\label{subsec:hmac}

Tanda tangan pada SeBatt dibuat dengan HMAC, yaitu fungsi yang menggabungkan kunci rahasia $K$
dengan sebuah pesan $m$ melalui fungsi hash \cite{rfc2104}. Fungsi hash yang dipakai adalah
SHA-256, yang menghasilkan keluaran 32 byte \cite{fips180}. HMAC dihitung dengan dua kali
pemanggilan hash:
\begin{equation}
  \mathrm{HMAC}(K, m) = H\big((K \oplus \mathit{opad}) \,\|\, H((K \oplus \mathit{ipad}) \,\|\, m)\big)
  \label{eq:hmac}
\end{equation}
dengan $H$ fungsi SHA-256, $\|$ penyambungan byte, serta $\mathit{ipad}$ dan $\mathit{opad}$
konstanta 64 byte bernilai \texttt{0x36} dan \texttt{0x5C} yang diulang. Keluaran HMAC disebut
\textit{tag}.

Ada dua sifat tag yang dipakai SeBatt. Pertama, perubahan satu bit saja pada pesan menghasilkan
tag yang sama sekali berbeda. Kedua, tag yang benar tidak dapat dibuat tanpa mengetahui $K$.
Kunci $K$ sepanjang 32 byte tersimpan di chip dan di basis data stasiun, dan tidak pernah
dikirim melalui saluran komunikasi.

Setiap sampel tegangan ditandatangani saat sampel diambil. Pesan yang ditandatangani terdiri
atas satu byte prefiks \texttt{0x53}, identitas baterai 4 byte, nomor urut 4 byte, dan tegangan
2 byte:
\begin{equation}
  \mathit{tag} = \mathrm{HMAC}\big(K,\; \texttt{0x53} \,\|\, \mathit{ID} \,\|\, \mathit{seq}
    \,\|\, V\big)
  \label{eq:tagcatatan}
\end{equation}
Stasiun menghitung ulang persamaan~\eqref{eq:tagcatatan} untuk setiap catatan yang diterimanya.
Apabila sebuah tegangan diubah setelah keluar dari chip, tag hasil hitungan stasiun tidak lagi
cocok dengan tag di catatan, sehingga riwayat ditolak dengan status \texttt{INVALID\_HMAC}.
Prefiks \texttt{0x53} membedakan pesan catatan dari dua pesan bertanda tangan lain pada
Bab~\ref{sec:arsitektur}, sehingga tag dari satu jenis pesan tidak dapat dipakai sebagai tag
jenis pesan lainnya.

\subsection{Nonce dan Serangan \textit{Replay}}
\label{subsec:nonce}

Tag pada persamaan~\eqref{eq:tagcatatan} membuktikan bahwa sebuah catatan berasal dari chip,
tetapi belum membuktikan bahwa catatan tersebut dikirim saat ini. Penyerang dapat merekam
balasan chip pada suatu waktu, lalu memutarnya ulang di kemudian hari, misalnya untuk menutupi
riwayat yang terjadi sesudah rekaman dibuat. Serangan semacam ini disebut \textit{replay}.

Untuk mencegahnya, stasiun membuat nonce, yaitu bilangan acak 32 byte yang baru untuk setiap
pemeriksaan. Nonce dikirim ke chip dan ikut ditandatangani di dalam balasan. Balasan hasil
rekaman memuat nonce lama, sehingga tag-nya tidak cocok dengan nonce baru yang sedang diperiksa
stasiun dan balasan tersebut ditolak dengan status \texttt{INVALID\_HMAC}. Stasiun juga
mencatat setiap nonce yang sudah dipakai. Pemeriksaan yang memakai nonce yang sama untuk kedua
kalinya ditolak dengan status \texttt{REPLAY\_DETECTED}.

\subsection{Nomor Urut Monoton}
\label{subsec:seq}

Tanda tangan per catatan belum cukup untuk mendeteksi catatan yang dihapus, karena catatan yang
tersisa tetap memiliki tag yang sah. Karena itu, setiap sampel diberi nomor urut (\textit{seq})
dari penghitung 32 bit di chip. Penghitung ini hanya bertambah satu setiap kali sampel diambil
dan tidak dapat diatur ulang melalui perintah apa pun. Akibatnya, riwayat yang utuh selalu
berupa deret seq yang berurutan tanpa celah.

Saat membaca riwayat, chip juga mengirim tag ringkasan yang mengikat nonce, identitas baterai,
seq tertua, dan seq terbaru di dalam buffer. Dengan tag ini stasiun mengetahui rentang seq yang
seharusnya diterima. Basis data stasiun menyimpan seq terakhir yang telah diverifikasi untuk
setiap baterai. Dari dua informasi tersebut, stasiun membedakan empat keadaan berikut.
\begin{itemize}
  \item \textbf{Deret utuh.} Seq berurutan dari seq tertua sampai seq terbaru dan setiap tag
    cocok. Riwayat diterima dengan status \texttt{VALID}.
  \item \textbf{Celah di tengah.} Ada seq yang hilang atau jumlah catatan tidak sesuai rentang
    yang ditandatangani. Riwayat ditolak dengan status \texttt{HISTORY\_GAP}.
  \item \textbf{Mundur.} Seq terbaru lebih kecil daripada seq terakhir di basis data, misalnya
    karena chip dikembalikan ke keadaan lama. Riwayat ditolak dengan status
    \texttt{HISTORY\_ROLLBACK}.
  \item \textbf{Tertimpa buffer.} Buffer chip hanya menampung 64 catatan, sehingga catatan
    tertua dapat tertimpa sebelum stasiun sempat membacanya. Selama sisa catatan berurutan dan
    seluruh tag cocok, riwayat tetap diterima dengan status \texttt{HISTORY\_TRUNCATED}.
\end{itemize}
Memotong catatan terbaru juga terdeteksi. Seq terbaru ikut ditandatangani di dalam tag
ringkasan, sehingga rentang yang dipersempit membuat tag ringkasan tidak cocok dan riwayat
ditolak dengan status \texttt{INVALID\_HMAC}.

\subsection{Pembacaan Tegangan Baterai}
\label{subsec:pembagi}

Tegangan baterai LiPo 3S berada di antara 9,0~V saat batas terkuras dan 12,6~V saat penuh,
jauh di atas batas masukan ADC sebesar 3,3~V. Tegangan tersebut diturunkan lebih dulu dengan
pembagi tegangan yang tersusun atas resistor $R_1$, potensiometer $P$, dan resistor $R_2$.
Dengan potensiometer di ujung atas, tegangan yang masuk ke ADC adalah
\begin{equation}
  V_{A0} = V_{\mathrm{bat}} \cdot \frac{P + R_2}{R_1 + P + R_2}
         = V_{\mathrm{bat}} \cdot \frac{10 + 15}{120 + 10 + 15}
         = \frac{V_{\mathrm{bat}}}{5{,}8}
  \label{eq:pembagi}
\end{equation}
sehingga tegangan 9,0--12,6~V menjadi 1,55--2,17~V di masukan ADC. Memutar potensiometer ke
bawah memperkecil rasio pembagi sampai $15/145$. Cara ini dipakai untuk memperagakan baterai
yang terkuras tanpa benar-benar menguras sel LiPo, sebagaimana dijelaskan pada
Bab~\ref{sec:rakit}.

Tegangan pada masukan A0 diubah menjadi bilangan 16 bit oleh ADS1115 \cite{ads1115}. Dengan
rentang $\pm$4,096~V, satu LSB setara dengan 0,125~mV, sehingga nilai dalam milivolt diperoleh
dengan menggeser hasil mentah tiga bit ke kanan:
\begin{equation}
  V_{A0}\,[\mathrm{mV}] = \frac{\mathit{raw}}{8}
  \label{eq:adc}
\end{equation}
FPGA membaca nilai tersebut melalui I2C dan menandatanganinya apa adanya. Nilai yang
ditandatangani adalah milivolt di pin A0, bukan tegangan baterai. Stasiun kemudian
mengalikannya dengan faktor 5,8 dari persamaan~\eqref{eq:pembagi} untuk menampilkan tegangan
baterai dan menilai kesehatannya. Dengan pembagian tugas ini, faktor tersebut dapat dikalibrasi
tanpa mengubah rancangan di FPGA.

% ============================================================
%  3. ARSITEKTUR DAN PROTOKOL
% ============================================================
\section{Arsitektur dan Protokol}
\label{sec:arsitektur}

\subsection{Komponen Sistem}

Purwarupa SeBatt terdiri atas tiga bagian, yaitu perangkat keras di sisi baterai, IP SeBatt
yang berjalan di FPGA, dan perangkat lunak stasiun di laptop. Gambar~\ref{fig:alurteknis}
memperlihatkan rantai ketiganya. Bagian A menunjukkan jalur data dari baterai sampai stasiun,
sedangkan bagian B menunjukkan urutan pesan saat baterai ditukar, yang dibahas pada
Subbab~\ref{subsec:alur}.

\gambarberkas{alur-teknis.png}{1.0}{Rantai perangkat dan urutan protokol SeBatt.}{fig:alurteknis}

Pada produk, IP SeBatt ditempatkan di dalam chip BMS baterai dan berkomunikasi dengan stasiun
melalui bus CAN. Purwarupa ini mengganti bus CAN dengan UART dan menjalankan stasiun beserta basis
datanya sebagai program Python di laptop. Protokol yang dipakai tetap sama, sehingga penggantian
saluran komunikasi tidak mengubah cara pemeriksaan riwayat.

\subsubsection*{Perangkat keras}

Tabel~\ref{tab:perangkat} merangkum perangkat keras yang dipakai beserta sambungannya. Semua
sinyal antara FPGA dan rangkaian luar melewati header Pmod J2 pada KV260. Nomor pin mengikuti
penomoran Pmod 12 pin, dengan pin 5 dan 11 sebagai GND serta pin 6 dan 12 sebagai 3,3~V.

\begin{table}[H]
  \centering
  \small
  \begin{tabularx}{\textwidth}{l X l}
    \toprule
    Komponen & Peran & Sambungan \\
    \midrule
    AMD Kria KV260 & Menjalankan IP SeBatt di logika terprogram (PL). Prosesor (PS) hanya
      menyediakan \textit{clock} 100~MHz dan sinyal reset & Pmod J2 \\
    LiPo 3S dan pembagi tegangan & Sumber tegangan yang dicatat. Kenop pada pembagi dipakai untuk
      memperagakan baterai terkuras & Ke A0 ADS1115 \\
    Modul ADS1115 & ADC 16 bit yang mengubah tegangan A0 menjadi milivolt & I2C, Pmod 1--2 \\
    Adaptor CP2102 & Saluran UART 115200~8N1 antara FPGA dan laptop & Pmod 7--8 \\
    LED hijau dan merah & Hasil RESULT terakhir dari stasiun, merah juga saat chip terkunci
      & Pmod 9 dan 10 \\
    LED \textit{heartbeat} & Berkedip sekitar 6~Hz sebagai tanda rancangan di PL berjalan & Pmod 4 \\
    Kabel ke GND & Memilih ADS1115 sebagai sumber sampel & Pmod 3 \\
    \bottomrule
  \end{tabularx}
  \caption{Perangkat keras purwarupa SeBatt.}
  \label{tab:perangkat}
\end{table}

Pin Pmod 3 memilih sumber sampel. Apabila pin ini terhubung ke GND, sampel diambil dari
ADS1115. Apabila terhubung ke 3,3~V, chip mencatat nilai tetap 3000~mV, yang berguna untuk
menguji jalur UART tanpa rangkaian ADC. Pin ini tidak boleh dibiarkan terbuka, karena pin yang
terbuka terbaca sebagai logika 1. Rangkaian lengkapnya dibahas pada Bab~\ref{sec:rakit}.

\subsubsection*{IP SeBatt di FPGA}

IP SeBatt ditulis dalam Verilog standar tanpa IP vendor dan tanpa prosesor lunak. Modul puncak
\texttt{kv260\_top} menerima \textit{clock} dari PS, menyelaraskan reset terhadap
\textit{clock}, dan mengeluarkan LED \textit{heartbeat}. Di bawahnya, \texttt{sebatt\_top}
menghubungkan dua bagian utama. Bagian pertama adalah \texttt{ads1115\_ctrl}, yang membaca ADC
melalui \texttt{i2c\_master}. Bagian kedua adalah \texttt{auth\_ctrl}, yang mengambil sampel
setiap 1~detik, memberi nomor urut, menandatangani sampel dengan \texttt{hmac\_sha256}, menyimpan
64 catatan terakhir, dan melayani perintah stasiun melalui \texttt{uart\_if}.

Jalur dari ADC sampai tanda tangan seluruhnya berupa rangkaian logika, tanpa perangkat lunak di
antaranya. Identitas baterai pada purwarupa ini bernilai \texttt{0x42415401}, yaitu teks
\texttt{BAT} diikuti byte \texttt{0x01}. Rincian setiap modul dibahas pada Bab~\ref{sec:rtl}.

\subsubsection*{Perangkat lunak stasiun}

Perangkat lunak stasiun terdiri atas dua program Python. Program utama berada di berkas
\texttt{model/station.py}. Program ini mengirim perintah ke chip, memeriksa setiap tanda tangan,
dan mengirim hasil pemeriksaan kembali ke chip. Basis datanya menyimpan empat hal untuk setiap
baterai, yaitu kunci $K$, seq terakhir yang telah diverifikasi, status kunci, dan nonce yang
sudah dipakai. Program ini juga menghitung kesehatan baterai dari riwayat yang lolos
pemeriksaan, seperti jumlah kejadian terkuras dalam dan siklus pengisian.

Program kedua adalah \textit{dashboard} web di folder \texttt{web/}, yang menampilkan hasil
pemeriksaan, grafik riwayat tegangan, dan kartu kesehatan baterai. \textit{Dashboard} memakai
aturan pemeriksaan dan penilaian kesehatan yang sama dengan \texttt{model/station.py}. Program ini
dapat berjalan dalam mode FPGA, yang terhubung ke board, atau mode simulasi, yang memakai model
chip di laptop untuk latihan tanpa board. Cara menjalankan kedua program dijelaskan pada
Bab~\ref{sec:panduan}.

\subsection{Format Frame dan Perintah}

\isi{Format [CMD][LEN big-endian][PAYLOAD][CRC-8 poly 0x31] dan tabel perintah 0x01--0x09
beserta isi payload. Prefiks tanda tangan: 0x41 autentikasi, 0x53 catatan, 0x48 ringkasan.}

\subsection{Alur Penukaran Baterai}
\label{subsec:alur}

\isi{Empat tahap: pencatatan sampel setiap 1 detik, AUTH\_REQ dengan nonce, READ\_HIST\_REQ
dengan tag ringkasan, lalu verifikasi di stasiun dan RESULT. Tiga RESULT gagal berturut-turut
mengunci chip.}

\subsection{Kode Hasil Pemeriksaan}

\begin{table}[H]
  \centering
  \small
  \begin{tabularx}{\textwidth}{l l X}
    \toprule
    Kode & Status & Arti \\
    \midrule
    0x00 & VALID & Tanda tangan cocok dan riwayat utuh \\
    0x01 & INVALID\_HMAC & Isi catatan atau balasan tidak cocok dengan tanda tangannya \\
    0x02 & REPLAY\_DETECTED & Nonce yang sama dipakai untuk kedua kalinya \\
    0x03 & BATTERY\_LOCKED & Chip terkunci setelah tiga kegagalan \\
    0x05 & HISTORY\_GAP & Ada nomor urut yang hilang di tengah riwayat \\
    0x06 & HISTORY\_ROLLBACK & Nomor urut chip lebih kecil dari catatan basis data \\
    0x07 & HISTORY\_TRUNCATED & Catatan lama tertimpa buffer 64, sisa riwayat sah \\
    \bottomrule
  \end{tabularx}
  \caption{Kode status yang dilaporkan stasiun.}
  \label{tab:status}
\end{table}

\isi{Satu paragraf: VALID dan HISTORY\_TRUNCATED dinyatakan DITERIMA, sisanya DITOLAK.}

% ============================================================
%  4. IMPLEMENTASI RTL
% ============================================================
\section{Implementasi RTL}
\label{sec:rtl}

\isi{Pengantar bab: hierarki kv260\_top, sebatt\_top, ads1115\_ctrl (i2c\_master), dan
auth\_ctrl (uart\_if, axil\_master, hmac\_sha256, sha256\_core). Semua modul Verilog standar
tanpa IP vendor dan tanpa prosesor lunak.}

\gambarplaceholder{Diagram hierarki modul RTL}

\subsection{Modul \texttt{hmac\_sha256}}
\isi{Antarmuka, empat blok SHA-256 per tanda tangan, 113 siklus per blok, potongan kode inti.}

\subsection{Modul \texttt{auth\_ctrl}}
\isi{Pewaktu sampel, penghitung seq, buffer 64 catatan, penanda tangan, parser frame dengan
P\_SKIP dan batas waktu RX, pembekuan buffer saat riwayat dikirim, penghitung gagal dan
penguncian, LED mengikuti RESULT terakhir.}

\subsection{Modul \texttt{ads1115\_ctrl} dan \texttt{i2c\_master}}
\isi{Mode single-shot, konfigurasi 0xC383, jeda 10 ms, cek bus macet, sinkronisasi SDA/SCL
dua flip-flop.}

\subsection{Modul \texttt{uart\_if} dan \texttt{kv260\_top}}
\isi{UART 115200 8N1. Top level: clock pl\_clk0 dari PS, sinkronisasi reset, LED heartbeat,
pemilih sumber src\_sel di Pmod 3, dan file constraint kv260.xdc.}

\subsection{Verifikasi}
\isi{Model referensi Python dan vektor uji, testbench cocotb beserta jumlah uji yang lolos,
hasil sintesis dan WNS. Hanya angka yang sudah dijalankan ulang.}

% ============================================================
%  5. PERAKITAN PERANGKAT
% ============================================================
\section{Perakitan Perangkat}
\label{sec:rakit}

\subsection{Daftar Perangkat dan Bahan}
\isi{Tabel: KV260 beserta adaptor 12 V dan kartu SD Ubuntu, kabel micro-USB, ADS1115, LiPo 3S,
R1 120 k$\Omega$, kenop 10 k$\Omega$, R2 15 k$\Omega$, kapasitor 100 nF, tiga LED dan resistor
330 $\Omega$, CP2102, kabel jumper, laptop dengan Vivado Hardware Manager dan Python 3.}

\subsection{Rangkaian Pembacaan Tegangan}
\gambarberkas{rangkaian-adc.png}{1.0}{Skematik rangkaian ADC, LED, dan sambungan Pmod J2.}{fig:skematik}
\isi{Langkah merakit bernomor dan tabel tegangan LiPo terhadap tegangan A0. Ukur wiper sebelum
disambung ke A0. Skematik sebaiknya dipindah ke halaman landscape atau dipecah dua.}

\subsection{Sambungan ke Header Pmod J2}
\isi{Tabel 12 pin beserta nama pin paket. Pmod 3 ke GND, pin 5V CP2102 tidak disambung.}

\subsection{Perangkat Lunak Stasiun}
\isi{Python 3, pyserial, dan uji tanpa board:}
\begin{lstlisting}[style=terminal]
py -3.14 -m model.station --offline --demo
\end{lstlisting}

% ============================================================
%  6. PANDUAN PENGGUNAAN
% ============================================================
\section{Panduan Penggunaan}
\label{sec:panduan}

\subsection{Memuat \textit{Bitstream} ke KV260}
\isi{Boot Ubuntu dari kartu SD, Vivado Hardware Manager, Auto Connect, xck26, Program Device.
LED heartbeat berkedip.}

\subsection{Pemeriksaan Awal}
\begin{lstlisting}[style=terminal]
py -3.14 -m model.station --check --port COM8
\end{lstlisting}
\isi{Contoh keluaran [OK] untuk STATUS, AUTH, READ\_HIST.}

\subsection{Menjalankan \textit{Dashboard}}
\begin{lstlisting}[style=terminal]
py -3.14 -m web.server --fpga COM8 --profil lipo3s --open
\end{lstlisting}
\gambarplaceholder{Tangkapan layar \textit{dashboard} dengan penanda nomor pada tiap bagian}
\isi{Arti tiap bagian dan warna pada \textit{dashboard}.}

\subsection{Skenario Penukaran Normal}
\isi{Langkah, hasil DITERIMA, LED hijau, kartu kesehatan terisi.}

\subsection{Skenario Baterai Terkuras Dalam}
\isi{Kenop diputar turun di bawah 9,0 V lalu kembali. Riwayat tetap DITERIMA dan kartu kesehatan
mencatat terkuras dalam.}

\subsection{Skenario Serangan}
\isi{Hapus catatan (HISTORY\_GAP), ubah angka (INVALID\_HMAC), muat ulang \textit{bitstream}
(HISTORY\_ROLLBACK). Untuk setiap serangan: langkah, hasil, LED, dan tangkapan layar.}

\subsection{Penguncian Chip dan Mode Terminal}
\isi{Tiga penolakan mengunci chip dan cara membukanya. Perintah interaktif station:}
\begin{lstlisting}[style=terminal]
py -3.14 -m model.station --port COM8 --interactive --profil lipo3s
\end{lstlisting}

\subsection{Kalibrasi Faktor Pembagi}
\isi{Kenop di ujung atas, ukur LiPo dengan multimeter, faktor = tegangan multimeter dibagi
pin mV, isi ke profil lipo3s.}

% ============================================================
%  7. HASIL UJI DI PERANGKAT KERAS
% ============================================================
\section{Hasil Uji di Perangkat Keras}
\label{sec:hasil}

\subsection{Susunan Uji}
\gambarplaceholder{Foto susunan uji: KV260, rangkaian LiPo dan ADS1115, CP2102, laptop}

\subsection{Kinerja Autentikasi dan Pembacaan Riwayat}
\isi{Tabel dari docs/catatan-teknis/hasil-uji-kv260.md: waktu AUTH rata-rata, jumlah
autentikasi per detik, waktu READ\_HIST 64 catatan, beserta tanggal uji dan bitstream yang
dipakai.} Nilai yang belum diukur ditulis sebagai \isian{belum diuji}.

\subsection{Deteksi Serangan}
\isi{Tabel setiap serangan dan status yang dihasilkan di board.}

% ============================================================
%  8. PEMECAHAN MASALAH DAN PERINGATAN
% ============================================================
\section{Pemecahan Masalah dan Peringatan}
\label{sec:masalah}

\subsection{Pemecahan Masalah}
\isi{Tabel gejala, penyebab, dan tindakan: port COM tidak muncul, tegangan selalu 3000 mV,
sampel 0 mV atau berhenti, HISTORY\_ROLLBACK setelah memuat ulang \textit{bitstream}, chip
terkunci.}

\subsection{Peringatan Penggunaan}
\isi{ADS1115 hanya 3,3 V. Kenop tidak boleh langsung ke 12 V. Jangan kuras LiPo di bawah
3,0 V per sel, dan cabut konektor balance setelah dipakai. Pmod 3 tidak dibiarkan terbuka.}

\subsection{Batasan Purwarupa}
\isi{ADC berupa modul eksternal, penghitung seq volatil, kunci simetris tetap dan tersimpan
di basis data, PUF belum terpasang, adegan terkuras dibuat dengan kenop. Pengembangan
berikutnya: kunci dari RO-PUF, penghitung di NVM, ADC satu die, sintesis SKY130.}

% ============================================================
%  DAFTAR PUSTAKA
% ============================================================
\clearpage
\phantomsection
\addcontentsline{toc}{section}{Daftar Pustaka}
% Urutan \bibitem mengikuti urutan sitasi pertama di teks, sehingga nomor [1] adalah sumber
% yang pertama disitasi. Sumber baru disisipkan sesuai posisi sitasi pertamanya.
\begin{thebibliography}{99}
  \bibitem{okezone} Okezone, ``Pengguna Kendaraan Listrik di Indonesia Terus Bertambah,''
    15 Oktober 2026. \url{https://digital.okezone.com/view/2026/10/15/4/106303/pengguna-kendaraan-listrik-di-indonesia-terus-bertambah}
  \bibitem{swapid} SWAP Energi, situs resmi. \url{https://www.swap.id/}
  \bibitem{electrum} Electrum, ``Tentang Kami.'' \url{https://electrum.id/id/about-us}
  \bibitem{detikvgreen} detikOto, ``Begini Sistem dan Biaya Sewa Baterai Motor Listrik
    VinFast,'' 2026. \url{https://oto.detik.com/motor-listrik/d-8581711/begini-sistem-dan-biaya-sewa-baterai-motor-listrik-vinfast}
  \bibitem{rukn} Kementerian Energi dan Sumber Daya Mineral, \textit{Keputusan Menteri ESDM
    Nomor 85.K/TL.01/MEM.L/2025 tentang Rencana Umum Ketenagalistrikan Nasional}, 5 Maret 2025.
    Ringkasan angka: CNBC Indonesia, ``Rencana Umum Ketenagalistrikan Resmi Dirilis, PLTU Batu
    Bara Masih Ada,'' 25 Maret 2025.
    \url{https://www.cnbcindonesia.com/news/20250325131108-4-621534/rencana-umum-ketenagalistrikan-resmi-dirilis-pltu-batu-bara-masih-ada}
  \bibitem{paparanruptl} Kementerian Energi dan Sumber Daya Mineral, \textit{Materi Konferensi
    Pers RUPTL PT PLN (Persero) Tahun 2025--2034: Meningkatkan Ketahanan dan Keberlanjutan Energi
    Nasional}, Jakarta, 26 Mei 2025.
    \url{https://gatrik.esdm.go.id/gatrik/gatrik-api/storage/migrasi/download_index/files/4ec39-materi-paparan-ruptl-2025-2034.pdf}
  \bibitem{databoksruptl} Databoks Katadata, ``Projected Number of Electric Vehicles in
    Indonesia 2025--2034,'' 3 Juli 2025, merangkum RUPTL PT PLN (Persero) 2025--2034.
    \url{https://databoks.katadata.co.id/en/transportation-logistics/statistics/68665063883b9/projected-number-of-electric-vehicles-in-indonesia-2025-2034}
  \bibitem{bq40z50} Texas Instruments, \textit{BQ40Z50-R5 Technical Reference Manual},
    SLUUCN4B, bagian 13.3--13.5. \url{https://www.ti.com/lit/ug/sluucn4b/sluucn4b.pdf}
  \bibitem{rfc2104} H.~Krawczyk, M.~Bellare, dan R.~Canetti, \textit{HMAC: Keyed-Hashing for
    Message Authentication}, RFC 2104, IETF, 1997.
  \bibitem{fips180} National Institute of Standards and Technology, \textit{Secure Hash Standard
    (SHS)}, FIPS PUB 180-4, 2015.
  \bibitem{ads1115} Texas Instruments, \textit{ADS111x Ultra-Small, Low-Power, I2C-Compatible,
    860-SPS, 16-Bit ADCs}, lembar data SBAS444.
  % Belum disitasi. Rencana dipakai di Bab 6 (memuat bitstream ke KV260).
  \bibitem{kv260} AMD, \textit{Kria KV260 Vision AI Starter Kit User Guide}, UG1089.
\end{thebibliography}

% ============================================================
%  LAMPIRAN
% ============================================================
\appendix
\renewcommand{\thesection}{\Alph{section}}

\section{Daftar Istilah}
\label{app:istilah}

\begin{table}[H]
  \centering
  \small
  \begin{tabularx}{\textwidth}{l X}
    \toprule
    Istilah & Arti \\
    \midrule
    ADC & \textit{Analog-to-Digital Converter}, pengubah tegangan menjadi angka \\
    \textit{Bitstream} & Berkas konfigurasi yang dimuat ke FPGA \\
    BMS & \textit{Battery Management System} \\
    FPGA & \isian{isi} \\
    HMAC & \isian{isi} \\
    I2C & \isian{isi} \\
    KV260 & \isian{isi} \\
    LiPo 3S & \isian{isi} \\
    Nonce & \isian{isi} \\
    Pmod & \isian{isi} \\
    PUF & \isian{isi} \\
    RTL & \isian{isi} \\
    Seq & \isian{isi} \\
    Stasiun penukaran (BSS) & \isian{isi} \\
    UART & \isian{isi} \\
    \bottomrule
  \end{tabularx}
  \caption{Daftar istilah.}
  \label{tab:istilah}
\end{table}

\section{Daftar Port Modul}
\label{app:port}

\isi{Tabel port untuk kv260\_top, sebatt\_top, auth\_ctrl, hmac\_sha256, ads1115\_ctrl,
i2c\_master, dan uart\_if: nama, arah, lebar, dan fungsi.}

\section{Kode Sumber}
\label{app:kode}

\isi{Kode lengkap dimuat dengan \texttt{\textbackslash lstinputlisting} dari folder rtl/ dan
fpga/kv260/ agar selalu sama dengan repositori. Urutan: kv260\_top.v, sebatt\_top.v,
auth\_ctrl.v, uart\_if.v, axil\_master.v, hmac\_sha256.v, sha256\_core.v, ads1115\_ctrl.v,
i2c\_master.v, kv260.xdc. Modul key\_vault.v dan ro\_puf.v untuk tahap 2 tidak termasuk.}

\subsection{Modul Puncak}
\subsection{Modul Autentikasi dan Komunikasi}
\subsection{Modul Kriptografi}
\subsection{Modul Pembacaan ADC}
\subsection{Constraint KV260}

\end{document}
